\section{KẾT LUẬN VÀ KIẾN NGHỊ}
\subsection{Khái quát hóa (generalization) về những điều tìm thấy}
    Đề tài đã hoàn thành xuất sắc các yêu cầu của Bài 10: Xây dựng thành công thuật toán tính toán lực tương tác tĩnh điện bằng phương pháp đại số ký hiệu (Symbolic Math) trên mặt phẳng Oxy.\par
\subsection{Ý nghĩa (implication)}
    Khẳng định hiệu quả vượt trội của việc kết hợp giữa lý thuyết Vật lý đại cương và các công cụ tin học hiện đại trong việc giải quyết bài toán lý thuyết phức tạp.\par
\subsection{Những hạn chế}
    \begin{itemize}[leftmargin=1.2cm]
    \item Bài toán mới xét trong không gian 2D ($Oxy$) và môi trường chân không liên tục.
    \end{itemize}
    \begin{itemize}[leftmargin=1.2cm]
    \item Chưa xét đến hiệu ứng phân cực trong các môi trường điện môi không đồng nhất.
    \end{itemize}
\subsection{Kiến nghị hướng nghiên cứu tiếp theo}
    \begin{itemize}[leftmargin=1.2cm]
    \item Phát triển chương trình lên không gian 3 chiều ($Oxyz$).
    \end{itemize}
    \begin{itemize}[leftmargin=1.2cm]
    \item Xây dựng giao diện người dùng GUI tương tác trực quan (Tkinter / PyQt).
    \end{itemize}